\section{Compactness of analytic interactions and coherent backgrounds}
\label{aic-section}

This is compactness of analytic local interactions and classical
maps at fixed parameters. It is not compactness of probability
measures or a continuum construction. It supplies a joint
infinite-volume interface for the total-kernel Ward and spatial
transport arguments.

\subsection{A limiting interaction at the original radius}
\label{aic-interaction-section}

Fix a cell size \(J\), radius \(b>0\), exponent \(a>0\), and
finitely many scalar coordinates per cell. Let compatible cubic
tori exhaust the lattice. Suppose their connected analytic
coefficients satisfy
\begin{equation}
 \sup_I\sum_{H\ni I}e^{a|H|}\|E_{N,H}\|_b\le\epsilon .
 \label{aic-input}
\end{equation}
No consistency between volumes, or symmetry of individual
coefficients, is assumed. For compactly supported fields define
the vacuum-normalized total
\begin{equation}
 \Phi_N(A)=\sum_H\bigl[E_{N,H}(A|_H)-E_{N,H}(0)\bigr].
 \label{aic-total}
\end{equation}
If \(K\) contains the cells meeting the field support, then
\(|\Phi_N(A)|\le2|K|\epsilon\) for \(\|A\|_\infty<b\).

\begin{proposition}[Interaction compactness without radius loss]
\label{aic-interaction-proposition}
There is a subsequence and a connected analytic family \(E_H\)
on the original open radius-\(b\) polydiscs with
\begin{equation}
 \sup_I\sum_{H\ni I}e^{a|H|}\|E_H\|_b\le\epsilon .
 \label{aic-limit-norm}
\end{equation}
Its vacuum-normalized total \(\Phi\) is absolutely convergent on
compact fields. On every fixed finite field support and every
radius \(r<b\), \(\Phi_N\) converges locally uniformly to \(\Phi\).
In particular all fixed derivatives of the totals converge.
\end{proposition}

\begin{proof}
Each fixed finite connected \(H\) has an unambiguous central
embedding on sufficiently large tori. Its coefficients are bounded
holomorphic functions on one finite-dimensional polydisc. Montel,
or diagonal extraction of bounded Taylor coefficients followed by
Cauchy estimates, gives a locally uniform subsequential limit.
There are countably many carriers, translates and scalar Taylor
indices, so one diagonal subsequence works for all of them.

For a finite selection of carriers, Fatou at radius \(r<b\)
passes their weighted bound. Increase \(r\) to \(b\), then
exhaust the carriers, to obtain \eqref{aic-limit-norm}.
The norm is the supremum on the open \(b\) polydisc. This
retains its original radius without asserting uniform convergence
at the boundary.

For carriers meeting \(K\), the omitted large-carrier terms obey
\begin{equation}
 \sum_{\substack{H\cap K\ne\varnothing\\|H|>Q}}
 |E_{N,H}(A|_H)-E_{N,H}(0)|
       \le2|K|\epsilon e^{-aQ}.
 \label{aic-carrier-tail}
\end{equation}
Only finitely many connected carriers of size at most \(Q\)
meet \(K\). Once the torus side is much larger than \(JQ\)
and the radius of \(K\), each has a unique central lift;
winding carriers are covered by the tail. Convergence on the
finite collection followed by \eqref{aic-carrier-tail} proves
convergence of the totals. Finite-dimensional Cauchy then
proves convergence of their fixed derivatives.
\end{proof}

The infinite kernels can therefore be derivatives of one analytic
interaction, rather than unrelated subsequential tensors.
Translations, signed permutations and reality of the totals pass
by applying the convergence to a compact field and its fixed
transform. For gauge invariance choose a compactly supported
small site gauge \(v\) such that \(A,A^v\) lie in a smaller
polydisc. The transformed field has support in the original
support plus a finite edge neighborhood. Passing the identity
\(\Phi_N(A^v)=\Phi_N(A)\) gives the limiting identity and
its differentiated Ward equations. Constant conjugation preserves
compact field support and passes as well. These are identities
of totals; no symmetry of \(E_H\) is asserted.

\subsection{Joint compactness of coherent complete-field families}
\label{aic-background-section}

Fix a preceding homogeneous depth \(M\) and endpoint ratio \(L\).
The current output spacing in final units is \(h=L^{-1}\).
Suppose the finite-regulator classical maps have the coherent
families of the supported pre-gauge construction:
\begin{align}
 a_N|_i&=\sum_{X\ne\varnothing}a_{N;i,X}(w_X),
       \qquad a_{N;i,X}(0)=0,\notag\\
 \sup_i\sum_X e^{\kappa\tau(\{i\}\cup X)+A|X|}
       \|a_{N;i,X}\|_{C^{1,\alpha};b}
       &\le Cb/L .
 \label{aic-background-norm}
\end{align}
The complete physical collar norm, including its maxima and
first-gradient H\"older quotients, is inside the support sum.
Coefficients with the same \(X\) restrict one global assignment
and agree on overlapping collars.

At fixed \(L\), each bounded collar has finitely many current
coordinates. Apply Montel and a countable diagonal extraction
to all fixed \(i,X\), output coordinates and Taylor indices.
Overlap identities pass to the limits. Each complete collar
norm is a finite collection of components and difference
quotients with fixed physical denominators. It is lower
semicontinuous under the convergence. Fatou thus preserves
\eqref{aic-background-norm} with the same constants.

The sum of terms with either
\(\tau(\{i\}\cup X)>Q\) or \(|X|>Q\) is bounded by
\[
 C(b/L)\bigl(e^{-\kappa Q}+e^{-AQ}\bigr).
\]
The remaining supports form a finite set, since they lie in
a fixed ball about \(i\). Hence the full complete fields
converge locally uniformly at each finite set of roots.
The limiting support series is normally convergent and
defines an analytic map \(a\). Uniform Cauchy remainders
also give its Banach-valued analyticity.

For compact final support \(K\), any nonzero coefficient has
\(X\cap K\ne\varnothing\). Its rooted tree therefore joins
\(i\) to \(K\), yielding
\begin{equation}
 \|a(w)\|_{\text{complete collar }i}
        \le C(b/L)e^{-\kappa\operatorname{dist}(i,K)}.
 \label{aic-background-tail}
\end{equation}
A bound summed over cells of \(K\) differs only by that finite
factor. The same estimate holds uniformly on centered finite
tori with periodic distance. It gives summable output tails
at fixed \(L\), including its fixed number of current coordinates
per unit cell.

This construction may be made jointly with the interaction
extraction and with any countable collection of fixed \(M,L\).
A cofinal sequence of tori can be chosen divisible by the first
finitely many required powers at each diagonal stage. The
result supplies coherent infinite families; it does not choose
a unique background independently of the subsequence.

\subsection{Summable fields and actual scalar composition}
\label{aic-composition-section}

The limiting interaction extends to summable coordinate fields
with \(\|A\|_\infty\le\rho<b\). Coordinate Cauchy and the
segment \(tA\) give
\[
 |E_H(A|_H)-E_H(0)|
 \le\frac{C\|E_H\|_b}{b-\rho}
                         \sum_{j\in H}|A_j|.
\]
Summing first over carriers containing each coordinate yields
\begin{equation}
 |\Phi(A)|\le\frac{C\epsilon}{b-\rho}\|A\|_{\ell^1},
 \qquad
 |\Phi(A)-\Phi(B)|
       \le\frac{C\epsilon}{b-\rho}\|A-B\|_{\ell^1}
 \label{aic-interaction-Lipschitz}
\end{equation}
when both fields lie in that interior ball. Coordinate counts
are already in the \(\ell^1\) sum; there is no volume factor.
The same estimates hold on the finite regulators.

Choose \(L\) so the classical amplitude bound places \(a_N(w)\)
in a fixed interior radius \(\rho<b\) on the final data domain.
For compact \(w\), \eqref{aic-background-tail} gives uniform
\(\ell^1\) tails of \(a_N(w)\). On centered representatives,
the periodic distance from an output root to \(K\) is at least
its centered distance from the origin minus the fixed radius
of \(K\), so this tail bound is uniform in volume.
Truncate the output to a large fixed region. On the region,
both the interaction coefficients and the classical fields
converge locally uniformly. Outside it use
\eqref{aic-interaction-Lipschitz}. It follows that
\begin{equation}
 \Phi_N(a_N(w))\longrightarrow\Phi(a(w))
 \label{aic-substitution-limit}
\end{equation}
locally uniformly on compact final supports and smaller data
radii. This is convergence of actual scalar compositions,
not only of formal Taylor series.

For decaying noncompact intermediate inputs, apply coordinate
Cauchy directly to \eqref{aic-background-norm}. Only supports
containing the input coordinate \(j\) contribute, and their
rooted tree spans \(i,j\). On an interior ball,
\begin{equation}
 \|\partial_{w_j}a_i\|
     \le\frac{C}{b-\rho}
                   e^{-\kappa\operatorname{dist}(i,j)}.
 \label{aic-map-derivative}
\end{equation}
The fixed input and output fiber dimensions, including \(L^4\),
are absorbed in \(C\). Summing over output roots proves an
\(\ell^1\)-to-\(\ell^1\) derivative bound, uniform in volume
at fixed \(L\). Segment integration gives the corresponding
Lipschitz estimate. Since \(a(0)=0\), summable inputs have
summable outputs.

The same truncation proof therefore applies to finite
compositions, provided every intermediate field remains in
a strict common interior of its analytic domain. Finite-volume
scalar identities, including a classical semigroup identity,
can pass to a jointly extracted finite list of maps and
interactions on compatible nested volumes. Separate arbitrary
subsequences need not compose.

\subsection{Scope}
\label{aic-scope}

The result jointly attains a subsequential analytic interaction
and coherent classical families at fixed parameters. It supplies
the total-kernel and complete-field hypotheses used by the
spatial remainder theorem. It proves neither uniqueness,
convergence of Gibbs measures, a continuum theory, nor consistency
of a quantum update for all input families. Descent of the actual
quantum map to a chosen interaction space, continuity there with
its normalizations and cutoffs retained, and infinite iteration
remain separate requirements.
